\documentclass[10pt,a4paper]{amsart}
\usepackage[margin=19mm]{geometry}
\usepackage{amsmath,amssymb,mathtools}
\usepackage[scr=rsfs,cal=euler]{mathalfa}
\usepackage{xfrac}
\usepackage[unicode,hidelinks,bookmarks=false]{hyperref}
\newcommand{\ie}{i.e.\ }
\newcommand{\cf}{cf.\ }
\newcommand{\resp}{respectively\ }
\newcommand{\Subsection}{\subsection}
\providecommand{\nobreakdash}{\nobreak\hbox{-}}
\setlength{\emergencystretch}{2em}
\multlinegap0pt
\usepackage{amssymb}
\usepackage{dsfont}
\newtheorem{thm}{Theorem}
\newtheorem{prop}{Proposition}
\newtheorem{lem}{Lemma}
\newcommand{\EEE}{\mathbb{E}}
\newcommand{\GDS}{\mathrm{GDS}}
\newcommand{\GDSd}{\mathrm{GDS}_\mathrm{d}}
\newcommand{\Grin}{\mathrm{GRIN}}
\newcommand{\LEPSr}{\mathrm{LEPS}_\mathrm{r}}
\newcommand{\RG}{\mathcal{R}(G)}
\newcommand{\RRR}{\mathbb{R}}
\newcommand{\dd}{\mathrm{d}}
\newcommand{\one}{\mathds{1}}
\newcommand{\twocdots}{\,\cdot\,,\cdot\,}
\title{On equitable scoring functions and optimal forecasting behaviour}
\author{Robert J. Taggart, Nicholas Loveday}
\date{Source: arXiv:2608.06710v1 (CC BY 4.0)}
\begin{document}
\maketitle

\noindent\emph{Source and licence.} On equitable scoring functions and optimal forecasting behaviour. Version arXiv:2608.06710v1 (CC BY 4.0), distributed by arXiv under the Creative Commons Attribution 4.0 International licence, http://creativecommons.org/licenses/by/4.0/, as recorded in the arXiv licence metadata for this exact version and shown on the abstract page https://arxiv.org/abs/2608.06710v1. Full licence notice: \texttt{licenses/arxiv-2608.06710-CC-BY-4.0.txt}.\par\smallskip

\noindent\textit{Editorial scope.} Counted unit: the article's prop scores (source0.tex lines 407--429). The article's complete original proof of it is delivered with it (source0.tex lines 1108--1147, one \texttt{proof} environment of 40 lines, which the article places away from the statement). No mathematics is written, repaired or re-derived: the statement and the proof are the article's own lines, in the article's own notation, and no retained line is substituted or re-flowed.  Retained uncounted as the source-local context the proof needs: lines 339--344; lines 341--343; lines 371--373; lines 379--384; lines 1045--1051; lines 1091--1093.  Nothing was omitted from any retained span, so the package carries no omission record.  No retained line contains a citation, so the wrapper carries no bibliography and no bibliography entry.  No retained line contains a figure environment, an image-inclusion command or drawing code, so no image is delivered and the package declares no asset dependency.  Editorial formatting only, in addition: an AMS \texttt{amsart} preamble that restates the article's own declarations and macros used by the retained lines, namely the environments \texttt{thm}, \texttt{prop}, \texttt{lem} on separate counters, together with the standard packages those lines need.  The retained environments print in this document as Theorem 1, Proposition 1, Lemma 1 in order of appearance, whereas the article prints the same items with the numbering of its own full document.  The complete native source of the article as distributed in the arXiv e-print for this exact version is bundled unchanged as \texttt{sources/207026/source0.tex}.
\begin{thm}\label{thm:GDS is equitable}
Suppose that $G\in\mathcal{P}$ and that $\mathcal{P}_0$ is a class of distributions contained in $\mathcal{P}$ such that the expected score $\EEE_G \GDS(F,Y;G,\mu)$ is finite for all $F$ in $\mathcal{P}_0$. Then the generalised diagonal score $\GDS(\twocdots ; G,\mu):\mathcal{P}_0\times I \to[0,\infty]$ is equitable with respect to $G$, with equitability constant $c$ given by
\begin{equation}\label{eq:GDS equitability constant}
c=\int_{\RG}\alpha(1-\alpha)\,\dd\mu(\alpha).
\end{equation}
\end{thm}

\begin{equation}
c=\int_{\RG}\alpha(1-\alpha)\,\dd\mu(\alpha).
\end{equation}

\begin{equation}\label{eq:S1}
S_1(x,y)= (\one\{x>y\}-G(y))^2 - (\one\{x>y\}-G(x))^2,
\end{equation}

\begin{align}
\GDSd(x,y;G,\mu)
&= \GDS(F_x,y;G,\mu) \notag\\
&= \int_{\RG} S_{\alpha,G^{-1}(\alpha)}(x,y)\,\dd\mu(\alpha) \notag \\
&=  \int_{\RG}  (\one\{y<x\} - \alpha)(\one\{G^{-1}(\alpha) < x\} - \one\{G^{-1}(\alpha) < y\}) \,\dd\mu(\alpha) \label{eq:deterministic GDS}
\end{align}

\begin{lem}\label{lem:alpha in RG}
Suppose that $G\in\mathcal{P}$ and $\alpha\in\RG$
\begin{enumerate}
\item[(i)] Then $G(G^{-1}(\alpha))=\alpha$.
\item[(ii)] Suppose also that $G$ is strictly increasing on $\mathcal{D}(G)$. Then $G^{-1}(\alpha) < x$ if and only if $\alpha < G(x)$.
\end{enumerate}
\end{lem}

\begin{lem}\label{lem:finite expectations}
If $G\in\mathcal{P}$ and $0<G(x)<1$ for all $x$ in $\RRR$ then $0\leq \EEE_G[-\ln(G(Y))]\leq 1$ and $0\leq\EEE_G[-\ln(1-G(Y))]\leq 1$.
\end{lem}

\begin{prop}\label{prop:real valued equitable scores}
Suppose that $G\in\mathcal{P}$.
\begin{enumerate}
\item[(i)] Suppose that $G$ satisfies either assumption (G1) or (G2), and the measure $\mu$ has no atoms on $(0,1)$.\footnote{Intuitively, $\mu$ has no atoms if the inclusion or exclusion of endpoints in the interval $[a,b]$ of integration, where $[a,b]\subset(0,1)$, does not change the value of the integral.} Then
\begin{equation}\label{eq:GDSd reduced}
\GDSd(x,y;G,\mu) = \int_{G(y)}^{G(x)}(\one\{x>y\} - \alpha)\,\dd\mu(\alpha) \qquad\forall x,y\in I.
\end{equation}
In particular, the score depends only on the interval between $G(y)$ and $G(x)$ in probability space.
\item[(ii)] If $G$ satisfies either (G1) or (G2) then the scoring function $S_1:I \times I\to[0,1]$ given by (\ref{eq:S1}) is equal to $\GDSd(\twocdots;G,\mu)$ when $\dd\mu(\alpha)=2\,\dd\alpha$, is equitable with respect to $G$, and is equivalent to $\LEPSr$. Its equitability constant is $1/3$ if $G$ satisfies (G1) and $1/3-G(0)^2+2G(0)^3/3$ if  $G$ satisfies (G2).
\item[(iii)] If $G$ satisfies either (G3) or (G4) then the scoring function $S_2: I\times I\to[0,\infty)$, defined by
\begin{equation}\label{eq:S2}
S_2(x,y)=\ln\left|\frac{\one\{x < y\} - G(x)}{\one\{x < y\} - G(y)}\right|, \qquad \forall x,y\in I,
\end{equation}
is equal to $\GDSd(\twocdots;G,\mu)$ when $\dd\mu(\alpha) = (\alpha(1-\alpha))^{-1}\,\dd\alpha$, is equitable with respect to $G$, and is equivalent to $\Grin$. The equitability constant for $S_2$ is $1$ under assumption (G3) and $1-G(0)$ under (G4).
\item[(iv)] Suppose that $u:(0,1)\to\RRR$ is bounded, nonnegative and measurable.\footnote{Measurability is a necessary condition for using $u$ in the integrand. Most functions that practitioners are familiar with are measurable.} Define functions $U$ and $V$ by
\[U(z)=\int_0^z u(\alpha)\,\dd\alpha \quad\text{and}\quad V(z)=\int_0^z \alpha u(\alpha)\,\dd\alpha\]
whenever $0\leq z \leq 1$. If $G$ satisfies either (G1) or (G2) then the scoring function $S_3:I \times I\to[0,\infty)$, defined by
\begin{equation}
S_3(x,y;u)= \one\{x>y\}(U(G(x)) - U(G(y))) - V(G(x)) + V(G(y)), \qquad \forall x,y\in I,
\end{equation}
is equal to $\GDSd(\twocdots;G,\mu)$ when $\dd\mu(\alpha)=u(\alpha)\,\dd\alpha$ and is equitable with respect to $G$. The equitability constant for $S_3$ is $\int_b^1 \alpha(1-\alpha)u(\alpha)\,\dd\alpha$, where $b=0$ under assumption (G1) and $b=G(0)$ under (G2).
\end{enumerate}
\end{prop} 

\begin{proof}[Proof of Proposition~\ref{prop:real valued equitable scores}]
To prove (i), suppose that $G$ satisfies either assumption (G1) or (G2) and the measure $\mu$ has no atoms. The continuity assumptions give $\RG=(0,1)$ under (G1) and $\RG=[G(0),1)$ under (G2). The strict monotonicity assumption on $G$ and Lemma~\ref{lem:alpha in RG} (i) gives
\begin{align*}
\{\alpha\in\RG:x\leq G^{-1}(\alpha)<y\}
&=\{\alpha\in\RG:G(x)\leq \alpha< G(y)\} \\
&=[G(x),G(y))
\end{align*}
whenever $x,y \in I$ and $x<y$. Then \eqref{eq:deterministic GDS} gives
\[
\GDSd(x,y;G,\mu)
=
\begin{cases}
\displaystyle\int_{[G(x),G(y))}  \alpha \,\dd\mu(\alpha), & x<y, \\
\displaystyle\int_{[G(y),G(x))} (1-\alpha) \,\dd\mu(\alpha), & x>y, \\
0, & x=y.
\end{cases}
\]
Equation~\eqref{eq:GDSd reduced} follows from the hypothesis that $\mu$ has no atoms.

To prove (ii), let $\dd\mu(\alpha)=2\,\dd\alpha$ and suppose that $G$ satisfies (G1) or (G2). Then by part (i),
\[
\GDSd(x,y;G,\mu)= \int_{G(y)}^{G(x)}2(\one\{x>y\} - \alpha)\,\dd\alpha = S_1(x,y).
\]
To prove equitability, note that $S_1(x,y)$ is bounded by $1$ and consequently the expected score $\EEE_G S_1(x,Y)$ is also bounded, so Theorem~\ref{thm:GDS is equitable} applies, and the equitability constant is readily calculated from \eqref{eq:GDS equitability constant} using $\RG=(0,1)$ under assumption (G1) or  $\RG=[G(0),1)$ under (G2). The equivalence of $S_1$ and $\LEPSr$ follows from the identity
\begin{equation}\label{eq:S1 LEPS equivalence}
S_1(x,y)=-\tfrac{1}{3}\,\LEPSr(x,y) + 2G(y)^2-2G(y)+\tfrac{2}{3}.
\end{equation}

To prove (iii), let $\dd\mu(\alpha)=(\alpha(1-\alpha))^{-1}\,\dd\alpha$ and suppose that $G$ satisfies (G3) or (G4). Then by part (i),
\[
\GDSd(x,y;G,\mu) = \int_{G(y)}^{G(x)} \frac{\dd\alpha}{\one\{x>y\} - \alpha} = S_2(x,y).
\]
To establish equitability, fix $x$ in $\RRR$ and note that
\begin{equation}\label{eq:S2 decomp}
\EEE_G S_2(x,Y)=\EEE_G \ln|\one\{x< Y\} - G(x)| + \EEE_G[-\one\{x< Y\}\ln(1-G(Y))] + \EEE_G[-\one\{x \geq Y\}\ln(G(Y))].
\end{equation}
By Lemma~\ref{lem:finite expectations} each term on the right hand side is finite, and hence by Theorem~\ref{thm:GDS is equitable} $S_2$ is equitable with respect to $G$. The equitability constant follows from \eqref{eq:GDS equitability constant} with $\RG=(0,1)$ under (G3) and $\RG=[G(0),1)$ under (G4). The identity $S_2(x,y)= -\Grin(x,y) - \ln(G(y)) - \ln(1-G(y)) - 1$ establishes that $S_2$ and $\Grin$ are equivalent.

The proof of (iv) proceeds similarly to (ii), noting that $\mu$ has no atoms and $\EEE_G S_3(x,Y)<\infty$.
\end{proof}

\end{document}
